% Part: computability
% Chapter: recursive-functions
% Section: primes

\documentclass[../../../include/open-logic-section]{subfiles}

\begin{document}

\olfileid{cmp}{rec}{pri}
\olsection{اوليه عددونه}

محدوده کميت ټاکنه او محدوده کمينه‌موندنه د دې دپاره ډېر وسايل راکوي چې وښيو طبيعي تابعې او اړيکې بنسټيزې بازګشتي دي۔ د بېلګې په توګه اړيکه «$x$، $y$ وېشي» په پام کښې ونيسئ، چې په~$x \mid y$ ليکل کېږي۔ اړيکه~$x \mid y$ هله صدق کوي چې~$y$ پر~$x$ بې له پاتې شوني ووېشل شي، يعنې~$y$ د~$x$ صحيح مضرب وي۔ که دا اړيکه صدق ونۀ کړي، يعنې کله چې د~$y$ پر~$x$ وېش د پاتې شوني په ژبه بيانېدے شي او پاتې شونی~$> 0$ وي، نو~$x \nmid y$ ليکو۔ په بله وينا، $x \mid y$ هله او يوازې هله چې د کوم~$z$ دپاره~$x \cdot z = y$ وي۔ ښکاره ده چې هر داسې~$z$، که موجود وي، بايد~$\leq y$ وي۔ نو~$x \mid y$ هله او يوازې هله چې د کوم~$z \le y$ دپاره~$x \cdot z = y$ وي۔ له~$=$ او ضرب څخه په محدودې وجودي کميت ټاکنې اړيکه~$x \mid y$ داسې تعريفولے شو:
\[
x \mid y \defiff \bexists{z \leq y}{(x \cdot z) = y}.
\]
نو مو وښودله چې~$x \mid y$ بنسټيزه بازګشتي ده۔
\textbf{د ژباړې د نسخې سمون (OLCMP-008):} سرچينې د وېش د اړيکې د ناکامېدو په تشريح کښې مقسوم او مقسوم‌عليه سرچپه کړي وو، او د صفر وېشونکي دپاره ئې هم د پاتې شوني ژبه کارولې وه۔ دلته لوری سم شوے او عمومي حالت ورپسې وجودي شرط ټاکي۔

طبيعي عدد~$x$ هله
\emph{اوليه} دے چې نۀ~$0$ وي، نۀ~$1$، او يوازې پر~$1$ او پر خپل ځان وېشل کېږي۔ په بله وينا، اوليه عددونه داسې دي چې کله هم~$y \mid x$ وي، نو يا~$y = 1$ وي يا~$y=x$۔ د دې د ازمېيلو دپاره چې~$x$ اوليه دے که نۀ، يوازې د ټولو~$y \le x$ دپاره $y \mid x$ ازمېيل پکار دي، ځکه که~$y > x$ وي، نو په خپله $y \nmid x$ وي۔ نو اړيکه~$\fn{Prime}(x)$، چې هله او يوازې هله صدق کوي چې~$x$ اوليه وي، داسې تعريفېدے شي:
\[
\fn{Prime}(x) \defiff x \geq 2 \land \bforall{y \leq x}{(y \mid x \lif y = 1 \lor y = x)}
\]
او ځکه بنسټيزه بازګشتي ده۔

اوليه عددونه~$2$، $3$، $5$، $7$، $11$، او داسې نور دي۔ تابع~$p(x)$ په پام کښې ونيسئ چې په دې لړۍ کښې~$x$م اوليه عدد راګرځوي؛ يعنې $p(0) = 2$، $p(1) = 3$، $p(2) = 5$، او داسې نور۔ (د اسانتيا دپاره به ډېر وخت~$p(x)$ د~$p_x$ په توګه ليکو: $p_0=2$، $p_1=3$، او داسې نور۔)

که يوه تابع~$\fn{nextPrime}(x)$ مو لرله چې له~$x$ څخه لومړنے لوے اوليه عدد راګرځوي، نو~$p$ په بنسټيز بازګښت په اسانۍ تعريفېدے شي:
\begin{align*}
  p(0) & = 2\\
  p(x+1) & = \fn{nextPrime}(p(x))
\end{align*}
څرنګه چې~$\fn{nextPrime}(x)$ هغه تر ټولو کوچنے~$y$ دے چې~$y > x$ وي او~$y$ اوليه وي، دا په بې‌پولې پلټنه په اسانۍ محاسبه کېدے شي۔ خو د اقليدس د يوې پايلې له امله ئې په محدودې کمينه‌موندنې هم تعريفولے شو: تل داسې يو اوليه عدد شته چې له~$x$ څخه لوے او له $\fact{x}+1$ څخه کوچنے يا برابر وي۔
\[
  \fn{nextPrime}(x) =
  \bmin{y \leq \fact{x}+1}{(y > x \land \fn{Prime}(y))}.
\]
دا ښيي چې~$\fn{nextPrime}(x)$، او ځکه~$p(x)$، يوازې محاسبه کېدونکې نۀ بلکې بنسټيزه بازګشتي ده۔
\textbf{د ژباړې د نسخې سمون (OLCMP-009):} سرچينې د راتلونکي اوليه عدد د تابع آرګومېنټ په دوو ځايونو کښې د تابع د نوم دننه ايښی و، حال دا چې نورې کارونې نوم او آرګومېنټ جلا ليکي۔ دلته دواړه ځایونه له هماغې رسمي کارونې سره يو شان شول۔

(که لېواله يئ، دا د اقليدس د قضيې يو لنډ ثبوت دے۔ که ورکړل شوے عدد له دوو کوچنے وي، 2 نېغ په نېغه مطلوب اوليه عدد دے۔ په پاتې حالت کښې فرض کړئ~$p_n$ تر ټولو لوے اوليه عدد دے چې~$\le x$ وي۔ د ټولو هغو اوليه عددونو حاصل‌ضرب $p = p_0\cdot p_1 \cdot \dots \cdot p_n$ په پام کښې ونيسئ چې $\le x$ دي۔ يا~$p+1$ اوليه دے، يا د~$x$ او~$p+1$ تر منځ بل اوليه عدد شته۔ ولې؟ فرض کړئ~$p+1$ اوليه نۀ دے۔ بيا کوم اوليه عدد $q \mid p+1$ شته چې~$q < p+1$ وي۔ اوليه عددونه~$\le x$ يو هم~$p+1$ نۀ وېشي۔ (د~$p$ د تعريف له مخې هر اوليه $p_i \le x$، $p$ وېشي، يعنې پاتې شونی~$0$ ورکوي۔ نو هر $p_i \le x$ د~$p+1$ په وېش کښې پاتې شونی~$1$ ورکوي، او ځکه $p_i \nmid p+1$۔) نو~$q$ يو اوليه عدد دے چې~$>x$ او~$<p+1$ دے۔ او $p \le \fact{x}$، نو داسې اوليه عدد شته چې~$>x$ او $\le \fact{x}+1$ وي۔)
\textbf{د ژباړې د نسخې سمون (OLCMP-010):} د سرچينې لنډ ثبوت سمدستي تر ورکړل شوي عدد پورې تر ټولو لوے اوليه عدد ټاکه، خو دا څيز د لومړنيو دوو طبيعي عددونو دپاره نشته۔ دلته واړه حالتونه جلا وارزول شول او د حاصل‌ضرب استدلال پاتې حالتونو ته محدود شو؛ د قضيې حد او د راتلونکي اوليه عدد د تابع تعريف نۀ دي بدل شوي۔

\begin{prob}
صحيح وېش~$d(x, y)$ د محدودې کمينه‌موندنې په کارولو تعريف کړئ۔
\end{prob}

\end{document}
